\section{分布式微分器的证明}

\subsection{预备知识}

\begin{assumption} \label{assump initially_connected}
    假设 $\mathcal{G}(0)$ 为无向且连通的图。其边权满足对每条初始边有 $a_{ij}(0)=a_{ji}(0)>0$，否则 $a_{ij}(0)=0$。
    此外，$\mathcal{B}=\operatorname{diag}(b_1,\ldots,b_N)\neq \mathbf{0}$，且 $b_i\ge 0$。
\end{assumption}

\begin{lemma}
    \label{lemma M>0}
    在假设 \ref{assump initially_connected} 下，矩阵 $M(0)=\mathcal{L}(0)+\mathcal{B}$ 正定。
\end{lemma}
\begin{proof}
    由于 $\mathcal{G}(0)$ 无向，$\mathcal{L}(0)$ 为对称半正定矩阵。
    由于图连通，$\operatorname{ker}\mathcal{L}(0)=\operatorname{span}\{\mathbf{1}_N\}$。
    对任意非零 $x\in\mathbb{R}^N$，
    \[
        x^\mathsf{T}M(0)x=x^\mathsf{T}\mathcal{L}(0)x+\sum_{i=1}^N b_i x_i^2\ge 0.
    \]
    若等号成立，则 $x=c\mathbf{1}_N$ 且对每个 $i$ 有 $b_i x_i^2=0$。
    由于至少有一个 $b_i$ 为正，故 $c=0$，与 $x\neq 0$ 矛盾。
    因此 $M(0)$ 正定。
\end{proof}

在证明引理 \ref{lemma estimator} 之前，先给出若干有用引理。

\begin{lemma}[有限时间比较]\label{lemma:finite_time}
    设 $V(t)\ge 0$ 绝对连续。若存在常数 $c>0$ 与 $\alpha\in(0,1)$，使得只要 $V(t)>0$ 就有
    \begin{equation}
        \dot{V}(t)\leq -cV(t)^\alpha.
    \end{equation}
    则 $V(t)$ 在有限时间内到达零。此外，收敛时间满足
    \begin{equation}
        T\leq \frac{V(0)^{1-\alpha}}{c(1-\alpha)}.
    \end{equation}
\end{lemma}
\begin{proof}
    在任意使 $V(t)>0$ 的区间上，对 $V^{1-\alpha}$ 求导得
    \[
        \frac{\mathrm{d}}{\mathrm{d}t}V(t)^{1-\alpha}
        =(1-\alpha)V(t)^{-\alpha}\dot V(t)
        \le -c(1-\alpha).
    \]
    故在 $V$ 到达零之前有 $V(t)^{1-\alpha}\le V(0)^{1-\alpha}-c(1-\alpha)t$，由此得到上述上界。
\end{proof}

\begin{lemma}[杨氏不等式]\label{lemma:generalized_young}
    设 $p,q>1$ 且 $1/p+1/q=1$。对任意 $x,\xi\in\mathbb{R}$，
    \begin{equation}
        |x\xi|\leq \frac{1}{p}|x|^p+\frac{1}{q}|\xi|^q.
    \end{equation}
\end{lemma}

\subsection{主体部分}

\begin{lemma}[ \cite{chenDistributedFinitetimeDifferentiator2024,aldana-lopezExactLeaderEstimation2025,sunDistributedFinitetimeFormation2025}]
    \label{lemma estimator}
    在假设 \ref{assump initially_connected} 下，下列系统在有限时间内收敛。
    该系统为超螺旋算法的分布式版本：
    \begin{equation}
        \begin{cases}
        \dot{e}_i=z_i-k_a \sig{\xi_i}^{1/2}\\
        \dot{z}_i=-k_b \sgn(\xi_i) + n_i(t)
    \end{cases},\quad i=1,\ldots,N
    \label{eq sta}
    \end{equation}
    其中 $e_i,z_i,\xi_i, n_i\in \mathbb{R}^2$，且 $\xi_i=\sum_{j\in\mathbb{N}_i(0)}a_{ij}(0)(e_i-e_j)+b_i e_i$。
    对于不连续的滑模系统，收敛性结论按 Filippov 意义理解。
    即 $\sgn(x)=x/\|x\|$（当 $\|x\|\neq 0$），而在 $x=\mathbf{0}$ 处表示闭单位球中的任意向量。
    我们定义 $\sig{x}^q=\|x\|^q \sgn(x)$。
    扰动 $n(t)$ 满足 $\sup_{t\geq 0, i\in \mathbb{N}}\|n_i(t)\| \leq \bar{n}$。
    参数 $k_a$ 与 $k_b$ 按下述规则选取：
    \begin{equation*}
      k_b >0,\quad k_b \geq \frac{\bar{n}}{\rho}, \quad 
      k_a \geq \sqrt{\frac{2\gamma_1+2\gamma_0+1}{\lambda_1}}\,\sqrt{k_b} > 0,
    \end{equation*}
    其中 $0<\rho<1$，$\gamma_0=(1 + 3\gamma_1)/(1 - \rho)$，$\gamma_1=1 + \rho$，而 $\lambda_1$ 为正定矩阵 $M(0)=\mathcal{L}(0)+\mathcal{B}$ 的最小特征值。
\end{lemma}

\begin{proof}
    定义堆叠向量
    \begin{align*}
        \xi &= [\xi_1^\mathsf{T}, \dots, \xi_N^\mathsf{T}]^\mathsf{T},\quad 
        e = [e_1^\mathsf{T}, \dots, e_N^\mathsf{T}]^\mathsf{T},\\
        \sig{\xi}^q &=[(\sig{\xi_1}^q)^\mathsf{T}, \dots, (\sig{\xi_N}^q)^\mathsf{T}]^\mathsf{T},\quad 
        n = [n_1^\mathsf{T}, \dots, n_N^\mathsf{T}]^\mathsf{T}.
    \end{align*}
    由图拉普拉斯矩阵的定义，
    \begin{equation}\label{eq:2_error_relation}
        \xi = ((\mathcal{L}(0)+\mathcal{B})\otimes I_2)e = (M(0)\otimes I_2) e.
    \end{equation}
    由假设 \ref{assump initially_connected} 与引理 \ref{lemma M>0}，$M(0)$ 对称且正定。
    引入缩放变量 $\nu_i = z_i/k_a$，并记 $\nu = [\nu_1^\mathsf{T},\dots,\nu_N^\mathsf{T}]^\mathsf{T}$。则动力学 \eqref{eq sta} 化为
    \begin{equation}
        \dot{e} = -k_a\bigl(\sig{\xi}^{1/2} - \nu\bigr),\qquad
        \dot{\nu} = -\frac{k_b}{k_a}\Bigl(\sgn(\xi) - \frac{n}{k_b}\Bigr). \label{eq:compact}
    \end{equation}

    构造 Lyapunov 函数候选 $V = V_1 + \gamma_0 V_2$，其中
    \begin{align*}
        V_1 &=\sum_{i\in\mathbb{N}}\Bigl(-\xi_i^\mathsf{T}\nu_i + \frac{1}{3}\|\nu_i\|^3 + \frac{2}{3}\|\xi_i\|^{3/2}\Bigr),\\
        V_2 &= \frac{1}{3}\sum_{i\in\mathbb{N}}\|\nu_i\|^3.
    \end{align*}
    对固定的 $\xi_i$，将 $V_1$ 的逐项关于 $\nu_i$ 极小化。记 $g(\nu_i)=-\xi_i^\mathsf{T}\nu_i+\frac13\|\nu_i\|^3$，其梯度为 $\|\nu_i\|\nu_i-\xi_i$，故驻点满足 $\|\nu_i\|\nu_i=\xi_i$，即 $\|\nu_i\|=\sqrt{\|\xi_i\|}$、$\nu_i=\xi_i/\sqrt{\|\xi_i\|}$，此时 $g(\nu_i)=-\frac23\|\xi_i\|^{3/2}$。由于 $\frac13\|\nu_i\|^3$ 为凸函数（其梯度 $\|\nu_i\|\nu_i$ 为单调算子），该驻点即全局最小，于是 $-\xi_i^\mathsf{T}\nu_i+\frac13\|\nu_i\|^3\ge -\frac23\|\xi_i\|^{3/2}$，从而 $-\xi_i^\mathsf{T}\nu_i+\frac13\|\nu_i\|^3+\frac23\|\xi_i\|^{3/2}\ge 0$，故 $V_1\ge 0$。
    由于 $\gamma_0>0$，$V=0$ 蕴含对所有 $i$ 有 $\nu_i=\mathbf{0}$，进而由 $V_1=0$ 得对所有 $i$ 有 $\xi_i=\mathbf{0}$。
    因此，$V$ 关于 $(\xi,\nu)$ 正定。

    \paragraph{$V$ 的上界。}
    将 $V_1$ 的每一项改写为
    \[
    -\xi_i^\mathsf{T}\nu_i + \|\nu_i\|^3 + \|\xi_i\|^{3/2} - \Bigl(\frac{2}{3}\|\nu_i\|^3 + \frac{1}{3}\|\xi_i\|^{3/2}\Bigr).
    \]
    对 $\|\nu_i\|^2\|\xi_i\|^{1/2}=\bigl(\|\nu_i\|^2\bigr)\bigl(\|\xi_i\|^{1/2}\bigr)$ 应用引理 \ref{lemma:generalized_young}（取 $p=3/2$、$q=3$），有 $\frac{2}{3}\|\nu_i\|^3 + \frac{1}{3}\|\xi_i\|^{3/2} \ge \|\nu_i\|^2\|\xi_i\|^{1/2}$。于是
    \begin{align*}
        -\xi_i^\mathsf{T}\nu_i + \frac{1}{3}\|\nu_i\|^3 + \frac{2}{3}\|\xi_i\|^{3/2}
        &\le -\xi_i^\mathsf{T}\nu_i + \|\nu_i\|^3 + \|\xi_i\|^{3/2} - \|\nu_i\|^2\|\xi_i\|^{1/2}\\
        &\le \bigl(\xi_i - \sig{\nu_i}^2\bigr)^\mathsf{T}\bigl(\sig{\xi_i}^{1/2}-\nu_i\bigr).
    \end{align*}
    利用 $\|\xi_i - \sig{\nu_i}^2\| \le \|\xi_i\| + \|\nu_i\|^2$ 以及估计
    \begin{align*}
        \|\xi_i\| + \|\nu_i\|^2 
        &\le \|\nu_i - \sig{\xi_i}^{1/2} - \nu_i\|^2 + \|\nu_i\|^2 \\
        &\le 3\|\nu_i\|^2 + 2\|\nu_i - \sig{\xi_i}^{1/2}\|^2,
    \end{align*}
    借助 Hölder 不等式与杨氏不等式可得
    \begin{align*}
        \bigl(\xi_i - \sig{\nu_i}^2\bigr)^\mathsf{T}
        \bigl(\sig{\xi_i}^{1/2}-\nu_i\bigr)
        &\le \bigl(3\|\nu_i\|^2
        + 2\|\nu_i - \sig{\xi_i}^{1/2}\|^2\bigr)
        \|\nu_i - \sig{\xi_i}^{1/2}\|\\
        &\le 2\|\nu_i\|^3 + 3\|\nu_i - \sig{\xi_i}^{1/2}\|^3.
    \end{align*}
    由此，
    \begin{equation}\label{eq:V_bound}
        V \le 3\sum_{i=1}^{N}\|\nu_i - \sig{\xi_i}^{1/2}\|^3 + \Bigl(2 + \frac{\gamma_0}{3}\Bigr)\sum_{i=1}^{N}\|\nu_i\|^3.
    \end{equation}

    \paragraph{$V$ 的时间导数。}
    沿 \eqref{eq:compact} 对 $V$ 求导得 $\dot{V}=\dot{V}_1 + \gamma_0 \dot{V}_2$，其中
    \begin{align*}
        \dot{V}_1 &=\sum_{i=1}^{N}\bigl(-\nu_i + \sig{\xi_i}^{1/2}\bigr)^\mathsf{T}\dot{\xi}_i 
                   + \sum_{i=1}^{N}\bigl(-\xi_i + \sig{\nu_i}^2\bigr)^\mathsf{T}\dot{\nu}_i,\\
        \dot{V}_2 &= \sum_{i=1}^{N}\sig{\nu_i}^2 \cdot \dot{\nu}_i.
    \end{align*}

    对 $\dot{V}_1$ 的第一部分，利用 $\xi = (M(0)\otimes I_2)e$、动力学 \eqref{eq:compact} 以及 Rayleigh–Ritz 定理（其中 $\lambda_1 = \lambda_{\min}(M(0))$），有
    \begin{align*}
        \sum_{i=1}^{N}\bigl(-\nu_i + \sig{\xi_i}^{1/2}\bigr)^\mathsf{T}\dot{\xi}_i
        &= -k_a\bigl(-\nu + \sig{\xi}^{1/2}\bigr)^\mathsf{T}(M(0)\otimes I_2)\bigl(-\nu + \sig{\xi}^{1/2}\bigr)\\
        &\le -k_a\lambda_1\sum_{i=1}^{N}\|\nu_i - \sig{\xi_i}^{1/2}\|^2.
    \end{align*}

    对 $\dot{V}_1$ 的第二部分，令 $w_i = -\xi_i + \sig{\nu_i}^2$。由 $\dot{\nu}_i = -\frac{k_b}{k_a}\sgn(\xi_i) + \frac{n_i}{k_a}$ 以及 $\|n_i\|\le \bar{n}\le k_b\rho$，得
    \begin{align*}
        w_i^\mathsf{T}\dot{\nu}_i 
        &= -\frac{k_b}{k_a}w_i^\mathsf{T}\sgn(\xi_i) + \frac{1}{k_a}w_i^\mathsf{T}n_i\\
        &\le \frac{k_b}{k_a}\|w_i\| + \frac{k_b\rho}{k_a}\|w_i\|
         = \frac{k_b}{k_a}\gamma_1\|w_i\|,
    \end{align*}
    其中 $\gamma_1 = 1+\rho$。利用前面得到的估计 $\|w_i\|\le 3\|\nu_i\|^2 + 2\|\nu_i - \sig{\xi_i}^{1/2}\|^2$，可得
    \begin{equation}
        \dot{V}_1 \le -\Bigl(k_a\lambda_1 - 2\gamma_1\frac{k_b}{k_a}\Bigr)\sum_{i=1}^{N}\|\nu_i - \sig{\xi_i}^{1/2}\|^2 
                  + 3\gamma_1\frac{k_b}{k_a}\sum_{i=1}^{N}\|\nu_i\|^2. \label{eq:dotV1}
    \end{equation}

    下面分析 $\dot{V}_2$：
    \begin{align*}
        \dot{V}_2 &= \frac{k_b}{k_a}\sum_{i=1}^{N}\sig{\nu_i}^2\cdot\Bigl(-\sgn(\xi_i) + \frac{n_i}{k_b}\Bigr)\\
        &= \frac{k_b}{k_a}\sum_{i=1}^{N}\sig{\nu_i}^2\cdot\Bigl(-\sgn(\nu_i) + \frac{n_i}{k_b} + \sgn(\nu_i) - \sgn(\xi_i)\Bigr).
    \end{align*}
    由于 $\|\sig{\nu_i}^2\| = \|\nu_i\|^2$ 且 $\sgn(\nu_i)^\mathsf{T}\frac{n_i}{k_b} \le \rho$，推得
    \begin{equation}
        \dot{V}_2 \le -\frac{k_b}{k_a}(1-\rho)\sum_{i=1}^{N}\|\nu_i\|^2 
                   + \frac{k_b}{k_a}\sum_{i=1}^{N}\sig{\nu_i}^2\cdot\bigl(\sgn(\nu_i) - \sgn(\xi_i)\bigr). \label{eq:dotV2_1}
    \end{equation}

    为估计交叉项，对每个 $i$ 定义
    \[
    f_i = 2\|\nu_i - \sig{\xi_i}^{1/2}\|^2 - \sig{\nu_i}^2\cdot\bigl(\sgn(\nu_i) - \sgn(\xi_i)\bigr).
    \]
    令 $x = \|\nu_i\|$，$y = \|\xi_i\|^{1/2}$，$c = \sgn(\nu_i)^\mathsf{T}\sgn(\xi_i)\in[-1,1]$。则
    \begin{align*}
        f_i &= 2(x^2 + y^2 - 2x y c) - (x^2 - x^2 c)\\
        &= (1+c)x^2 + 2y^2 - 4cxy.
    \end{align*}
    我们证明 $f_i\ge 0$。当 $c\ge 0$ 时，$f_i$ 是关于 $x$ 的二次式：$f_i(x) = (1+c)x^2 - 4c y x + 2y^2$，其最小值在 $x^* = \frac{2c}{1+c}y$ 处取得，最小值为 $\frac{2(1+c-2c^2)}{1+c}y^2 = \frac{2(1-c)(1+2c)}{1+c}y^2 \ge 0$。当 $c<0$ 时，对称轴为负，故在 $x\ge 0$ 上最小值在 $x=0$ 处，此时 $f_i(0)=2y^2\ge 0$。因此恒有 $f_i\ge 0$，从而
    \begin{equation}
        \sig{\nu_i}^2\cdot\bigl(\sgn(\nu_i) - \sgn(\xi_i)\bigr) \le 2\|\nu_i - \sig{\xi_i}^{1/2}\|^2. \label{eq:cross_bound}
    \end{equation}
    将 \eqref{eq:cross_bound} 代入 \eqref{eq:dotV2_1} 得
    \begin{equation}
        \dot{V}_2 \le -\frac{k_b}{k_a}(1-\rho)\sum_{i=1}^{N}\|\nu_i\|^2 
                   + 2\frac{k_b}{k_a}\sum_{i=1}^{N}\|\nu_i - \sig{\xi_i}^{1/2}\|^2. \label{eq:dotV2_final}
    \end{equation}

    \paragraph{$\dot V$ 的综合。}
    将 \eqref{eq:dotV1} 与 \eqref{eq:dotV2_final} 结合，利用 $\dot V = \dot V_1 + \gamma_0 \dot V_2$：
    \begin{align*}
        \dot V \le &-\Bigl(k_a\lambda_1 - 2\gamma_1\frac{k_b}{k_a} - 2\gamma_0\frac{k_b}{k_a}\Bigr)\sum_{i=1}^{N}\|\nu_i - \sig{\xi_i}^{1/2}\|^2 \\
        &+ \Bigl(3\gamma_1\frac{k_b}{k_a} - \gamma_0\frac{k_b}{k_a}(1-\rho)\Bigr)\sum_{i=1}^{N}\|\nu_i\|^2.
    \end{align*}
    取 $\gamma_0 = \frac{1+3\gamma_1}{1-\rho}$。则 $\|\nu_i\|^2$ 的系数变为 $\frac{k_b}{k_a}[3\gamma_1 - \gamma_0(1-\rho)] = -\frac{k_b}{k_a}$。
    条件 $k_a \ge \sqrt{\frac{2\gamma_1+2\gamma_0+1}{\lambda_1}}\sqrt{k_b}$ 保证
    \[
    k_a\lambda_1 \ge \bigl(2\gamma_1+2\gamma_0+1\bigr)\frac{k_b}{k_a},
    \]
    这使得 $\|\nu_i - \sig{\xi_i}^{1/2}\|^2$ 的系数不大于 $-\frac{k_b}{k_a}$。于是
    \begin{equation}\label{eq:V_dot}
        \dot V \le -\frac{k_b}{k_a}\Bigl(\sum_{i=1}^{N}\|\nu_i - \sig{\xi_i}^{1/2}\|^2 + \sum_{i=1}^{N}\|\nu_i\|^2\Bigr).
    \end{equation}

    \paragraph{有限时间收敛。}
    令
    \[
        s_1 = \sum_{i=1}^{N}\|\nu_i - \sig{\xi_i}^{1/2}\|^3,
        \qquad
        s_2 = \sum_{i=1}^{N}\|\nu_i\|^3 .
    \]
    由 $L^p$ 范数关系（因 $3>2$，$\|\cdot\|_3\le \|\cdot\|_2$ 推出 $\sum a_i^3 \le (\sum a_i^2)^{3/2}$），有
    \[
    \sum_{i=1}^{N}\|\nu_i - \sig{\xi_i}^{1/2}\|^2 \ge s_1^{2/3},\qquad 
    \sum_{i=1}^{N}\|\nu_i\|^2 \ge s_2^{2/3}.
    \]
    不等式 \eqref{eq:V_bound} 可写为 $V \le 3 s_1 + (2+\frac{\gamma_0}{3}) s_2 \le C_{\max}(s_1+s_2)$，
    其中 $C_{\max} = \max\{3,\;2+\frac{\gamma_0}{3}\}$。
    利用 $(a+b)^{2/3} \le a^{2/3}+b^{2/3}$（$a,b\ge 0$），得
    \[
    \dot V \le -\frac{k_b}{k_a}\bigl(s_1^{2/3}+s_2^{2/3}\bigr)
           \le -\frac{k_b}{k_a}(s_1+s_2)^{2/3}
           \le -\frac{k_b}{k_a}\,C_{\max}^{-2/3}\,V^{2/3}.
    \]
    由引理 \ref{lemma:finite_time}（取 $\alpha=2/3$），$V$ 在有限时间内收敛到 $0$，且收敛时间满足
    \begin{equation}
        T \le 3\,\frac{k_a}{k_b}\,C_{\max}^{2/3}\,V(0)^{1/3}.
    \end{equation}
    当 $V=0$ 时，对所有 $i$ 有 $\nu_i = \xi_i = 0$，从而 $z_i=0$；又因 $M(0)$ 可逆，得 $e = (M(0)^{-1}\otimes I_2)\xi = 0$。证毕。
\end{proof}

\subsection{时变扰动上界的情形}

上面的证明只用到扰动满足 $\|n_i(t)\|\le \bar n$ 这一**逐时刻**的上界，全过程从未对 $\bar n$ 求导，也从未要求其为常数。因此，若扰动上界是时变的，即已知
\[
    \|n_i(t)\|\le \bar n(t),\qquad \bar n^*:=\sup_{t\ge 0}\bar n(t)<\infty,
\]
则只需在选取增益时改用其一致上界 $\bar n^*$：取 $\rho\in(0,1)$ 并令
\[
    k_b\ge \frac{\bar n^*}{\rho},
\]
其余参数（$k_a$、$\gamma_0$、$\gamma_1=1+\rho$、$\lambda_1$）的选取规则完全不变。

证明中关键的一步（对 $\dot V_1$ 的第二部分）有
\[
    w_i^\mathsf{T}\dot\nu_i
    =-\frac{k_b}{k_a}w_i^\mathsf{T}\sgn(\xi_i)+\frac{1}{k_a}w_i^\mathsf{T}n_i
    \le \frac{k_b}{k_a}\|w_i\|+\frac{\bar n(t)}{k_a}\|w_i\|
    \le \frac{k_b}{k_a}\|w_i\|+\frac{\bar n^*}{k_a}\|w_i\|
    =\frac{k_b}{k_a}(1+\rho)\|w_i\|,
\]
与 $\bar n$ 为常数的情形逐字相同——区别仅在于把常数 $\bar n$ 替换为一致上界 $\bar n^*$。由此
\[
    \dot V\le -\frac{k_b}{k_a}\Bigl(\sum_{i=1}^{N}\|\nu_i-\sig{\xi_i}^{1/2}\|^2+\sum_{i=1}^{N}\|\nu_i\|^2\Bigr)
\]
依然在每个时刻成立，从而系统仍在有限时间内**精确**收敛到原点，且收敛时间上界 $T\le 3(k_a/k_b)C_{\max}^{2/3}V(0)^{1/3}$ 保持不变。

结论：只要时变上界具备有限的上确界，原设计（定常增益）无需任何结构性修改即可适用；唯一的变化是增益选取条件中的 $\bar n$ 应理解为 $\sup_{t\ge0}\bar n(t)$。

\subsection{扰动上界无界的情形：时变增益设计}

若 $\bar n(t)$ 随时间无界增长（即 $\sup_{t\ge0}\bar n(t)=\infty$），任何固定的 $k_b$ 都会在足够大的 $t$ 处被扰动超过，符号项 $-k_b\sgn(\xi_i)$ 不再能主导 $n_i(t)$。因此定常增益设计不再适用。下面给出一种时变增益的充分条件。

设扰动满足
\[
    \|n_i(t)\|\le \bar n(t),\qquad i=1,\ldots,N,
\]
其中 $\bar n(t)$ 可无界，但存在一个绝对连续的正函数 $K_b(t)$ 作为在线上包络，使
\begin{equation}\label{eq:unbounded_Kb}
    K_b(t)\rho\ge \bar n(t),\qquad
    0<K_{b0}\le K_b(t),\qquad t\ge0 .
\end{equation}
取
\begin{equation}\label{eq:time_varying_gain_design}
    k_b(t)=K_b(t),\qquad
    k_a(t)=c_a\sqrt{K_b(t)},
\end{equation}
其中
\[
    c_a>\sqrt{\frac{2\gamma_1+2\gamma_0+1}{\lambda_1}},\qquad
    \gamma_0=\frac{1+3\gamma_1}{1-\rho},\quad \gamma_1=1+\rho .
\]
由于 $k_a$ 现在随时间变化，缩放变量 $\nu_i=z_i/k_a(t)$ 的动力学多出一项：
\[
    \dot\nu
    =-\frac{k_b(t)}{k_a(t)}\Bigl(\sgn(\xi)-\frac{n(t)}{k_b(t)}\Bigr)
    -\frac{\dot k_a(t)}{k_a(t)}\nu .
\]
前一项的估计与定常增益情形相同；后一项需要额外约束增益变化速度。记
\[
    S(\xi,\nu):=\sum_{i=1}^{N}\|\nu_i-\sig{\xi_i}^{1/2}\|^2+\sum_{i=1}^{N}\|\nu_i\|^2 .
\]
对任意给定初值，由 $V$ 的正定性与径向无界性，可在初始子水平集
\[
    \Omega_{V(0)}:=\{(\xi,\nu):0\le V(\xi,\nu)\le V(0)\}
\]
上取常数 $C_T(V(0))>0$，使由 $-\frac{\dot k_a}{k_a}\nu$ 产生的附加项满足
\begin{equation}\label{eq:time_varying_extra_bound}
    \left|
    \frac{\dot k_a(t)}{k_a(t)}
    \sum_{i=1}^{N}\Bigl(\xi_i^\mathsf T\nu_i-(1+\gamma_0)\|\nu_i\|^3\Bigr)
    \right|
    \le
    C_T(V(0))\frac{|\dot k_a(t)|}{k_a(t)}S(\xi,\nu).
\end{equation}
这里之所以可以取有限的 $C_T(V(0))$，是因为 $\Omega_{V(0)}$ 为紧集，且当 $S(\xi,\nu)=0$ 时必有 $\xi=\nu=0$，附加项的分子也同时为零。
因此，只要存在 $\theta\in(0,1)$ 使
\begin{equation}\label{eq:slow_gain_condition}
    C_T(V(0))\frac{|\dot k_a(t)|}{k_a(t)}
    \le
    \theta\frac{k_b(t)}{k_a(t)},\qquad t\ge0,
\end{equation}
则定常增益证明中的负定项可以吸收时变缩放项。具体地，有
\[
    \dot V
    \le
    -(1-\theta)\frac{k_b(t)}{k_a(t)}S(\xi,\nu).
\]
又由 \eqref{eq:time_varying_gain_design} 与 \eqref{eq:unbounded_Kb}，
\[
    \frac{k_b(t)}{k_a(t)}
    =\frac{\sqrt{K_b(t)}}{c_a}
    \ge \frac{\sqrt{K_{b0}}}{c_a}>0.
\]
结合前文的 $V$ 上界 \eqref{eq:V_bound} 与 $L^p$ 范数关系，得到
\[
    \dot V
    \le
    -\frac{1-\theta}{c_a}\sqrt{K_{b0}}\,
    C_{\max}^{-2/3}V^{2/3}.
\]
由有限时间比较引理 \ref{lemma:finite_time}，系统仍在有限时间内精确收敛到原点，且
\[
    T
    \le
    3\,\frac{c_a}{(1-\theta)\sqrt{K_{b0}}}\,
    C_{\max}^{2/3}V(0)^{1/3}.
\]

由于 $k_a(t)=c_a\sqrt{K_b(t)}$，条件 \eqref{eq:slow_gain_condition} 可写成更直接的形式
\[
    \frac{1}{2}C_T(V(0))\left|\frac{\dot K_b(t)}{K_b(t)}\right|
    \le
    \theta\frac{\sqrt{K_b(t)}}{c_a}.
\]
也就是说，$K_b(t)$ 可以随无界扰动上界增长，但其相对增长速度不能过快。若 $\bar n(t)$ 有跳变或增长过快，应先构造一个绝对连续、满足 \eqref{eq:unbounded_Kb} 与上述慢变条件的光滑上包络 $K_b(t)$，再由 \eqref{eq:time_varying_gain_design} 生成 $k_a(t),k_b(t)$。

例如，当 $\bar n(t)$ 绝对连续时，可取
\[
    K_b(t)=\frac{\bar n(t)+\delta_b}{\rho},\qquad
    \delta_b>0,
\]
并令
\[
    k_b(t)=\frac{\bar n(t)+\delta_b}{\rho},\qquad
    k_a(t)=c_a\sqrt{\frac{\bar n(t)+\delta_b}{\rho}}.
\]
此时只需验证
\[
    \frac{C_T(V(0))}{2}
    \frac{|\dot{\bar n}(t)|}{\bar n(t)+\delta_b}
    \le
    \frac{\theta}{c_a}
    \sqrt{\frac{\bar n(t)+\delta_b}{\rho}},
\]
等价地，
\[
    |\dot{\bar n}(t)|
    \le
    \frac{2\theta}{c_a C_T(V(0))\sqrt{\rho}}\,
    \bigl(\bar n(t)+\delta_b\bigr)^{3/2}.
\]
